% Part: normal-modal-logic
% Chapter: frame-correspondence
% Section: standard-translation
%READERNOTE{SOL6-A153: मानक भाषांतराच्या प्रत्येक मोडल पायरीत नवीन बद्ध चल निवडण्याची अट स्पष्ट केली. Monadic हे द्वितीय-क्रम संख्यापकांपुरते बंधन आहे; व्यक्तिचलांवरील प्रथम-क्रम संख्यापक अनुमत आहेत. परावर्तितेच्या उदाहरणात व्यक्तिचलाचे मूल्य, प्रत्यक्ष जग आणि successor set वेगळे करून नेमणूक स्पष्ट केली.}

\documentclass[../../../include/open-logic-section]{subfiles}

\begin{document}

\olfileid{nml}{frd}{st}

\olsection{द्वितीय-क्रमातील परिभाष्यता}

मोडल सूत्रांनी परिभाषित होणारा चौकटींचा प्रत्येक गुणधर्म
प्रथम-क्रमात परिभाष्य असेलच असे नाही. पण एकस्थानी
विधेयांवर संख्यापन करण्याची परवानगी दिल्यास (म्हणजे
एकस्थानी द्वितीय-क्रम संख्यापन), मोडल सूत्रांनी परिभाष्य
असलेले सर्व चौकट-गुणधर्म परिभाषित करता येतात.
एखादे मोडल !!{formula} जगात सत्य असण्याच्या अटी
प्रथम-क्रम !!{formula}s शी ज्या पद्धतशीर रीतीने
संबंधित असतात, तिचा उपयोग करणे ही युक्ती आहे.
या पद्धतीला मोडल सूत्रांचे प्रथम-क्रम सूत्रांत
मानक भाषांतर म्हणतात. त्यासाठीच्या भाषेत
प्राप्यता संबंधासाठी द्विस्थानी !!{predicate}~$Q$
तर असतेच, शिवाय~$!A$ मध्ये येणाऱ्या
!!{propositional variable}s~$\Obj p_i$ साठी
एकस्थानी !!{predicate}~$P_i$ देखील असते.

\begin{defn}
  \emph{मानक भाषांतर}~$\ST_x(!A)$ ची विगमनाने
  पुढीलप्रमाणे व्याख्या करतो. मोडल उपवाक्यांत $x$ पेक्षा वेगळा
  नवीन व्यक्तिचल $y$ निवडा. आतील भाषांतरातील संख्यापकांसाठीही
  मुक्त चल बद्ध होणार नाहीत अशी नवीन बद्ध चलांची नावे निवडावीत:
  \begin{enumerate}
  \tagitem{prvFalse}{\indcase{!A}{\lfalse}{$\ST_x(\indfrm) = \lfalse$.}}{}
  \tagitem{prvTrue}{\indcase{!A}{\ltrue}{$\ST_x(\indfrm) = \ltrue$.}}{}
  \item \indcase{!A}{\Obj p_i}{$\ST_x(\indfrm) = \Atom{P_i}{x}$.}
  \tagitem{prvNot}{\indcase{!A}{\lnot !B}{$\ST_x(\indfrm) = \lnot \ST_x(!B)$.}}{}
  \tagitem{prvAnd}{\indcase{!A}{(!B \land !C)}{$\ST_x(\indfrm) = (\ST_x(!B)
      \land \ST_x(!C))$.}}{}
  \tagitem{prvOr}{\indcase{!A}{(!B \lor !C)}{$\ST_x(\indfrm) = (\ST_x(!B)
      \lor \ST_x(!C))$.}}{}
  \tagitem{prvIf}{\indcase{!A}{(!B \lif !C)}{$\ST_x(\indfrm) = (\ST_x(!B)
      \lif \ST_x(!C))$.}}{}
  \tagitem{prvIff}{\indcase{!A}{(!B \liff !C)}{$\ST_x(\indfrm) = (\ST_x(!B)
      \liff \ST_x(!C))$.}}{}
  \tagitem{prvBox}{\indcase{!A}{\Box !B}{$\ST_x(\indfrm) =
      \lforall[y][(\Atom{Q}{x,y} \lif \ST_y(!B))]$.}}{}
  \tagitem{prvDiamond}{\indcase{!A}{\Diamond !B}{$\ST_x(\indfrm) =
      \lexists[y][(\Atom{Q}{x,y} \land \ST_y(!B))]$.}}{}
  \end{enumerate}
\end{defn}

उदाहरणार्थ, $\ST_x(\Box p \lif p)$ हे
$\lforall[y][(\Atom{Q}{x,y}
  \lif \Atom{P}{y})] \lif \Atom{P}{x}$ असे आहे.
~$\ST_x(!A)$ च्या भाषेतील कोणत्याही !!{structure}
साठी प्रांत, $Q$ ला नेमून दिलेला द्विस्थानी संबंध,
आणि एकस्थानी !!{predicate}s~$P_i$ यांना नेमून
दिलेले प्रांताचे उपसंच लागतात. दुसऱ्या शब्दांत,
अशा !!a{structure} चे घटक हे~$!A$ च्या प्रतिमानाच्या
घटकांशी तंतोतंत जुळतात: प्रांत म्हणजे जगांचा संच;
$Q$ ला नेमलेला द्विस्थानी संबंध म्हणजे प्राप्यता
संबंध; आणि $P_i$ यांना नेमलेले उपसंच म्हणजे
नेमणुका~$V(\Obj p_i)$. त्यामुळे मोडल प्रतिमानातील
$!A$ ची पूर्ति आणि अनुरूप !!{structure} मधील
$\ST_x(!A)$ ची पूर्ति जुळते, हे अपेक्षितच आहे:

\begin{prop}\ollabel{prop:st}
  $\mModel{M} = \tuple{W, R, V}$ असू द्या.
  $\Domain{M'} = W$, $\Assign{Q}{M'} = R$, आणि
  $\Assign{P_i}{M'} = V(\Obj p_i)$ असलेली
  $\Struct{M'}$ ही प्रथम-क्रम !!{structure} असू द्या,
  तसेच $s(x) = w$ असू द्या. मग
  \[
  \mSat{M}{!A}[w] \text{ तेव्हा आणि केवळ तेव्हाच } \Sat{M'}{\ST_x(!A)}[s]
  \]
\end{prop}

\begin{proof}
  $!A$ वर विगमनाने.
\end{proof}

\begin{prop}
  $!A$ हे मोडल !!{formula} आणि
  $\mModel{F} = \tuple{W, R}$ ही चौकट असू द्या.
  $\Domain{F'} = W$ आणि $\Assign{Q}{F'} = R$
  असलेली $\Struct{F'}$ ही प्रथम-क्रम !!{structure}
  असू द्या. पुढील द्वितीय-क्रम सूत्राला $!A'$
  म्हणा:
  \[
  \lforall[X_1][\dots\lforall[X_n][\lforall[x][
        \SSubst{\ST_x(!A)}{\subst{X_1}{P_1}, \dots,
          \subst{X_n}{P_n}}]]],
  \]
  जिथे $P_1$, \dots, $P_n$ ही~$\ST_x(!A)$
  मधील सर्व एकस्थानी !!{predicate}s आहेत. मग
  \[
  \mModel{F} \Entails !A \text{ तेव्हा आणि केवळ तेव्हाच } \Sat{F'}{!A'}
  \]
\end{prop}

\begin{proof}
  $\Sat{F'}{!A'}$ तेव्हा आणि केवळ तेव्हाच, जेव्हा
  $i = 1$, \dots,~$n$ साठी
  $\Assign{P_i}{M'} \subseteq W$ असलेल्या
  मूळ रचनेच्या प्रत्येक विस्ताररूप
  !!{structure}~$\mModel{M'}$ मध्ये, आणि
  $s(x) \in W$ असलेल्या प्रत्येक $s$ साठी,
  $\Sat{M'}{\ST_x(!A)}[s]$ असते.
  \olref{prop:st} नुसार, हे तेव्हा आणि केवळ
  तेव्हाच, जेव्हा $\mModel{F}$ वर आधारित सर्व
  प्रतिमाने~$\mModel{M}$ आणि प्रत्येक जग~$w \in W$
  यांसाठी $\mSat{M}{!A}[w]$ असते; म्हणजे
  $\mModel{F} \Entails !A$.
\end{proof}

\begin{defn}
  चौकटींचा वर्ग~$\mClass{F}$ \emph{द्वितीय-क्रमात
  परिभाष्य} असतो, जर एकच द्विस्थानी
  !!{predicate}~$P$ असलेल्या आणि द्वितीय-क्रमाचे संख्यापक
  केवळ एकस्थानी संचचलांवर असलेल्या भाषेत असे
  !!a{sentence}~$!A$ असेल की,
  $\Domain{M} = W$ आणि $\Assign{P}{M} = R$
  असलेल्या !!{structure}~$\Struct{M}$ मध्ये
  $\Sat{M}{!A}$ तेव्हा आणि केवळ तेव्हाच,
  जेव्हा $\mModel{F} = \tuple{W, R} \in \mClass{F}$.
  व्यक्तिचलांवरील प्रथम-क्रम संख्यापकही अनुमत आहेत.
\end{defn}

\begin{cor}
  चौकटींचा वर्ग एखाद्या !!a{formula}~$!A$ ने
  परिभाष्य असेल, तर त्याच्या अनुरूप प्राप्यता
  संबंधांचा वर्ग एकस्थानी द्वितीय-क्रम वाक्याने
  परिभाष्य असतो.
\end{cor}

\begin{proof}
  आधीच्या सिद्धतेतील एकस्थानी द्वितीय-क्रम
  वाक्य~$!A'$ ला आवश्यक गुणधर्म आहे.
\end{proof}

उदाहरण म्हणून पुन्हा !!{formula}~$\Box p \lif p$
पाहा. ते परावर्तितेची व्याख्या करते. अर्थात,
परावर्तिता !!{sentence}~$\lforall[x][\Atom{Q}{x,x}]$
याने प्रथम-क्रमात परिभाष्य आहे. पण ती पुढील
एकस्थानी द्वितीय-क्रम !!{sentence} नेही
परिभाष्य आहे:
\[
\lforall[X][\lforall[x][(\lforall[y][(\Atom{Q}{x,y} \lif \Atom{X}{y})]
    \lif \Atom{X}{x})]].
\]
म्हणजे ही दोन वाक्ये समतुल्य आहेत. आधी द्वितीय-क्रम वाक्य
!!a{structure}~$\Struct{M}$ मध्ये सत्य आहे असे समजा.
कोणतेही $w \in W$ घ्या. $R=\Assign{Q}{M}$ असू द्या आणि
$S=\Setabs{z \in W}{Rwz}$ ठेवा. $s(x)=w$ आणि $s(X)=S$
असलेली नेमणूक $s$ निवडा. वैश्विक संख्यापनाच्या व्याख्येवरून
\[
\Sat{M}{\lforall[y][(\Atom{Q}{x,y} \lif \Atom{X}{y})]
    \lif \Atom{X}{x}}[s].
\]
$S$ च्या निवडीवरून या अभिव्यंजनाचे पूर्वांग $s$ नुसार सत्य आहे:
$w$ पासून प्राप्य असलेले प्रत्येक जग $S$ मध्ये आहे.
म्हणून $\Sat{M}{\Atom{X}{x}}[s]$; म्हणजे $w \in S$.
$S$ च्या व्याख्येनुसार $Rww$. $w$ कोणतेही असल्यामुळे
$\Sat{M}{\lforall[x][\Atom{Q}{x,x}]}$ मिळते.

आता $\Sat{M}{\lforall[x][\Atom{Q}{x,x}]}$
असे समजा आणि
$\Sat{M}{\lforall[X][\lforall[x][(\lforall[y][(\Atom{Q}{x,y} \lif
        \Atom{X}{y})] \lif \Atom{X}{x})]]}$
हे दाखवू. कोणतीही नेमणूक~$s$ घ्या आणि
$\Sat{M}{\lforall[y][(\Atom{Q}{x,y} \lif
    \Atom{X}{y})]}[s]$ असे माना. $s'$ ही~$s$
ची $y$-भिन्नरूप नेमणूक असू द्या, जिथे $s'(y)
= s(x)$; मग
$\Sat{M}{\Atom{Q}{x,y} \lif \Atom{X}{y}}[s']$,
म्हणजे $\Sat{M}{\Atom{Q}{x,x} \lif \Atom{X}{x}}[s]$.
$\Sat{M}{\lforall[x][\Atom{Q}{x,x}]}$ असल्यामुळे
पूर्वांग सत्य आहे; म्हणून $\Sat{M}{\Atom{X}{x}}[s]$,
आणि हेच दाखवायचे होते.

चौकटींचे काही परिभाष्य वर्ग प्रथम-क्रमात परिभाष्य
नसल्यामुळे, $!A'$ या रूपातील प्रत्येक एकस्थानी
द्वितीय-क्रम !!{sentence} हे प्रथम-क्रम
!!{sentence} शी समतुल्य असेलच असे नाही.
कोणती अशी आहेत हे ठरवण्याची प्रभावी पद्धत नाही.

\end{document}
